% table_hw.tex — Table IV: ESP32-P4 hardware benchmark.
% Average over three prompts; each run generates 64 tokens.

\begin{table}[t]
\centering
\footnotesize
\caption{On-device benchmark on ESP32-P4 Nano
(RISC-V HP, 360\,MHz, 32\,MB PSRAM).}
\label{tab:hw}
\setlength{\tabcolsep}{3pt}
\renewcommand{\arraystretch}{1.15}
\begin{tabular}{lrrrr}
\toprule
\textbf{Method} & \textbf{Model} & \textbf{PSRAM} &
\textbf{ms/tok} & \textbf{tok/s} \\
 & \textbf{KB} & \textbf{free KB} & & \\
\midrule
INT4U  & 7{,}431 & 21{,}566 & $1{,}040.6\pm0.6$ & $0.961$ \\
INT4BW & 9{,}285 & 19{,}682 & $1{,}355.8\pm0.4$ & $0.737$ \\
CTLE   & 7{,}433 & 21{,}566 & $\mathbf{1{,}035.9\pm1.1}$ & $\mathbf{0.965}$ \\
P-CTLE & 7{,}433 & 21{,}566 & $\mathbf{832.5\pm1.5}$ & $\mathbf{1.201}$ \\
\bottomrule
\end{tabular}
\smallskip
\begin{minipage}{0.98\columnwidth}
  \footnotesize
  Values are means over prompts with 4, 8, and 6 prefill tokens.
  INT4BW is slower because each row reads per-group scales from PSRAM.
  P-CTLE preserves the CTLE footprint and replaces the inner-loop multiply
  with a product-LUT lookup after 4-bit activation quantisation.
\end{minipage}
\end{table}
